% ==============================================================================
% Slide 07: Financial Decomposition & DuPont Analysis
% ==============================================================================
\begin{frame}{Financial Decomposition \& DuPont Analysis}{Deconstructing Equity Returns into Orthogonal Operational and Financial Drivers}
    \vspace{-0.20cm}
    \begin{columns}[T,onlytextwidth]
        % Column 1: DuPont Formulations
        \begin{column}{0.485\textwidth}
            \begin{mathcard}[height=2.74cm]{3-Stage DuPont Decomposition}
                \fontsize{5.7pt}{6.9pt}\selectfont
                Deconstructs headline ROE into three orthogonal business drivers:
                \begin{equation*}
                    \textstyle \text{ROE} = \underbrace{\frac{\text{Net Income}}{\text{Sales}}}_{\text{Margin (P1)}} \times \underbrace{\frac{\text{Sales}}{\text{Assets}}}_{\text{Efficiency (P5)}} \times \underbrace{\frac{\text{Assets}}{\text{Equity}}}_{\text{Leverage (P2)}}
                \end{equation*}
                \begin{itemize}
                    \setlength{\itemsep}{1.2pt}
                    \item \textbf{Diagnostic Insight:} Decouples whether superior ROE stems from pricing power, asset velocity, or risky balance sheet gearing.
                    \item \textbf{Driver Attribution:} Pinpoints margin vs.\ leverage trade-offs.
                \end{itemize}
            \end{mathcard}

            \vspace{0.12cm}

            \begin{mathcard}[height=2.74cm]{5-Stage Extended DuPont Identity}
                \fontsize{5.6pt}{6.8pt}\selectfont
                Separates financing and tax burdens from operating performance:
                \begin{equation*}
                    \textstyle \text{ROE} = \underbrace{\frac{\text{NI}}{\text{EBT}}}_{\text{Tax}} \times \underbrace{\frac{\text{EBT}}{\text{EBIT}}}_{\text{Interest}} \times \underbrace{\frac{\text{EBIT}}{\text{Sales}}}_{\text{Operating Margin}} \times \underbrace{\frac{\text{Sales}}{\text{Assets}}}_{\text{Turnover}} \times \underbrace{\frac{\text{Assets}}{\text{Equity}}}_{\text{Leverage}}
                \end{equation*}
                \begin{itemize}
                    \setlength{\itemsep}{1.2pt}
                    \item \textbf{Extended Guardrail:} Isolates transient tax incentives or debt refinancing relief from enduring operational margin contraction.
                    \item \textbf{Operational Integrity:} Validates true core profitability.
                \end{itemize}
            \end{mathcard}
        \end{column}

        % Column 2: Cross-Sector Insights & Production Utility
        \begin{column}{0.485\textwidth}
            \begin{quantcard}[height=2.74cm]{Cross-Sector Comparative Normalization}
                \fontsize{5.8pt}{7.0pt}\selectfont
                \begin{itemize}
                    \setlength{\itemsep}{1.5pt}
                    \item \textbf{IT Services:} Zero debt, moderate asset turnover, and high net margins ($20\text{--}25\%$).
                    \item \textbf{Utilities \& Infra:} High leverage ($\frac{\text{Assets}}{\text{Equity}} > 3.5$) with low asset turnover and narrow operating margins.
                    \item \textbf{Automobile \& FMCG:} Rapid inventory turnover offsetting competitive gross margins.
                \end{itemize}
            \end{quantcard}

            \vspace{0.12cm}

            \begin{quantcard}[height=2.74cm]{Analytical Risk Early-Warning}
                \fontsize{5.8pt}{7.0pt}\selectfont
                \begin{itemize}
                    \setlength{\itemsep}{1.5pt}
                    \item \textbf{Leverage Masking Flag:} Intercepts firms sustaining high ROE purely through escalating debt leverage as operating margins contract.
                    \item \textbf{Automated Alert:} Triggers \badge[jade]{Leverage Overhang} warning badge in the FinMetrics Suite rating engine.
                    \item \textbf{Executive Memo:} Directly feeds AI narration context.
                \end{itemize}
            \end{quantcard}
        \end{column}
    \end{columns}
\end{frame}
